\chapter{Tecnologias que Serão Utilizadas}
\label{cap:tecnologias}

% Descrição sucinta das tecnologias que podem ser utilizadas para desenvolver cada módulo do sistema:
% - Escolha justificada da tecnologia de implementação de cada módulo.
% - Descrever com mais detalhes as tecnologias escolhidas.
% - As tecnologias alternativas podem ser apresentadas de maneira mais resumida.

\section{Tecnologias em Avaliação}

As tecnologias de implementação ainda não foram definidas de forma
definitiva. Como o projeto combina análise estática, representações
estruturadas de código e modelos de \textit{Machine Learning}, a escolha
será realizada por meio de provas de conceito, considerando o suporte a
Java e C\#, o custo computacional, a integração entre os módulos e a
reprodutibilidade dos resultados.

\subsection{Linguagens de Entrada e de Implementação}

Java e C\# foram definidas como as linguagens de entrada da versão inicial
da ferramenta por serem utilizadas em ambientes corporativos. A linguagem
de implementação do \textit{pipeline}, entretanto, será selecionada durante
o desenvolvimento. Python é uma alternativa devido ao seu ecossistema de
ML; Java ou Kotlin podem facilitar a integração com a JVM; e C\#/.NET
permite o uso de APIs de compilação e análise de projetos dessa linguagem.

A decisão considerará a experiência da equipe e as bibliotecas disponíveis.
Formatos estruturados, como JSON, poderão ser utilizados para reduzir o
acoplamento entre os módulos.

\subsection{Análise Sintática, Estrutural e Estática}

A representação inicial poderá utilizar árvores sintáticas abstratas,
complementadas por informações de fluxo de controle e de dados. O uso de
\textit{Code Property Graphs} será avaliado por meio de uma prova de
conceito, pois os CPGs reúnem essas representações em um único grafo
\cite{yamaguchi2014cpg}.

CPG/Codyze e Joern serão considerados para a geração e consulta dos grafos
\cite{codyze2026,joern2026}. A escolha dependerá da qualidade dos grafos
produzidos para Java e C\# e do custo de integração. Se o CPG não apresentar
ganho suficiente, serão mantidas representações mais simples, como AST,
CFG e DFG. SonarQube, CodeQL, Semgrep e analisadores específicos das
linguagens poderão fornecer alertas complementares.

\subsection{Modelos de Classificação}

O primeiro modelo deverá ser um classificador de referência com baixo
custo computacional, utilizado para validar as características extraídas e
as métricas de avaliação. A biblioteca dependerá da linguagem escolhida;
em Python, o scikit-learn constitui uma alternativa para classificação e
validação cruzada \cite{scikitlearn2026}. Posteriormente, poderão ser avaliados CodeBERT e GraphCodeBERT
\cite{feng2020codebert,guo2021graphcodebert}. No ecossistema Python,
PyTorch e Transformers são alternativas para executar esses modelos
\cite{huggingface2026}. Eles não serão requisitos para o produto mínimo
viável devido ao maior custo computacional.

% \subsection{LLMs, Relatórios e Integração}

% O uso de \textit{Large Language Models} será complementar. O LLM poderá
% auxiliar na explicação e na priorização dos alertas, mas não será a única
% fonte para determinar uma vulnerabilidade. A interface deverá evitar
% dependência de um único modelo ou fornecedor.

% Os resultados serão apresentados inicialmente em linha de comando e em
% arquivos estruturados, contendo a CWE, a localização, a confiança da
% classificação e a justificativa do alerta, quando disponíveis.

\subsection{Recursos de \textit{Hardware}}

O projeto não prevê o desenvolvimento de \textit{hardware} específico. O
desenvolvimento e os testes iniciais serão realizados nos computadores
pessoais da equipe, utilizando CPU e memória convencionais. Caso o
treinamento ou a avaliação de modelos mais complexos exija maior capacidade,
poderão ser utilizadas GPUs locais ou infraestrutura em nuvem. A necessidade
será definida conforme o modelo e o volume de dados selecionados.

\subsection{Ambiente de Desenvolvimento e Execução}

O código e as configurações serão mantidos em um sistema de controle de
versão, o Git. As dependências serão registradas e, quando aplicável, Docker será
utilizado para isolar analisadores externos e padronizar a execução.

\section{Tecnologias Alternativas Avaliadas}

As alternativas serão comparadas por meio de provas de conceito com um
conjunto reduzido de arquivos Java e C\#. Serão considerados o suporte às
linguagens, o tempo de execução, a integração com os demais módulos e a
qualidade das características produzidas. Tecnologias mais complexas, como CPGs e modelos pré-treinados, somente
serão adotadas se apresentarem ganho mensurável sobre a linha de base. Caso
contrário, será mantida a alternativa mais simples para preservar a
viabilidade do cronograma e a entrega do protótipo.
